\section{Main Theorem and Proof Outline}

\begin{theorem}[Main theorem]\label{thm: main theorem}
    Let $\mathcal{F}$ be a complex signature set that is closed under conjugation. 
    Then $\Holant(\mathcal{F})$ is \#P-hard unless $\mathcal{F}$ satisfies {\rm (\ref{cond: tractable classes})}, in which case it is tractable.
\end{theorem}

\paragraph{Proof Overview}
The tractable cases of the main theorem follow from \cref{thm: tractable classes}.
It remains to prove that $\Holant(\mathcal{F})$ is \#P-hard if $\mathcal{F}$ does not satisfy condition~(\ref{cond: tractable classes}).

We first generalize the second order orthogonality ({\sc 2nd-Orth}) in \cite{realholant} to the conjugate closed setting.
This requires a genuinely new argument heavily depending on a recent dichotomy on quaternary signatures (\cref{thm: single quaternary dichotomy}).
Throughout the rest of the proof, {\sc 2nd-Orth} plays an important role.

Suppose $\mathcal{F}$ fails condition \eqref{cond: tractable classes}.
In particular, $\mathcal{F}\not\subseteq \mathscr{T}$.
Since $\mathcal{U}^{\otimes}\subseteq \mathscr{T}$, we also have $\mathcal{F}\not\subseteq \mathcal{U}^{\otimes}$.
Thus, there is a nonzero signature $f\in\mathcal{F}$ of arity $2n$ such that $f\notin\mathcal{U}^{\otimes}$.
We prove \#P-hardness by induction on $2n$.
The base case $2n=2$ is handled by \textsc{1st-Orth} (\cref{lm: not 1st-orth is hard}).
The case $2n=4$ is handled by \textsc{2nd-Orth} and the nonexistence of an $\operatorname{AME}(4,2)$ state (\cref{lm: base case 2n=4}).

For higher arities, we aim to realize a signature $f'$ from $f$ with lower arity while keeping the property $f'\notin\mathcal{U}^{\otimes}$.
Regular approaches for arity reduction are merging, mating, and factorization.
However, this approach does not always succeed due to the existence of the special signatures $f_6$ and $f_8$ discovered in~\cite{realholant}.
These signatures themselves are not in $\mathcal{U}^{\otimes}$, but the lower-arity signatures obtained from them by merging, mating, and factorization all lie in $\mathcal{U}^{\otimes}$.
We therefore need separate proofs for $2n=6$ and $2n=8$.

For $2n=6$, our idea is to first realize $f_6$, and then prove that $\Holant(f_6,\mathcal{F})$ is \#P-hard under the assumption that $\mathcal{F}\not\subseteq\mathscr{A}$.
In \cref{sec: Third Order Orth}, we first realize an $\operatorname{AME}(6,2)$ state, or equivalently, a signature satisfying \textsc{3rd-Orth}.
We use the framework of Xia~\cite{MingjiProgram}, which shows that, unless hardness already follows, the ``projective binary group'' generated by $f$ is a finite subgroup of $\mathbf{PU}(2)=\mathbf{SU}(2)/\{\pm I_2\}$.
We give an exposition of this proof in \autoref{appendix: mingji}.
Since $\mathbf{SU}(2)/\{\pm I_2\}\cong\mathbf{SO}(3)$, we use the classification of finite subgroups of $\mathbf{SO}(3)$ to consider all possible cases of the projective binary group in \cref{sec: Third Order Orth}.
For the polyhedral group case (\cref{subsec: polyhedral groups}), the large dihedral group case (\cref{subsec: dihedral group}), and the standard Klein group case (\cref{subsubsec: Standard Case Klein}), we prove that a signature satisfying \textsc{3rd-Orth} can be realized.
We also prove that the non-standard Klein group case (\cref{subsubsec: Non-standard Klein}) and the cyclic group cases (\cref{subsec: Trivial group,subsec: Order Two group,subsec: Cyclic group of Order at least three}) cannot occur.


Following the realization of an $\operatorname{AME}(6,2)$ state in~\cref{sec: Third Order Orth}, we show how to recover $f_6$ in the~\cref{sec:recover-f6}. By the uniqueness of $\operatorname{AME}(6,2)$ states up to local
unitary transformations and input permutations, which is proved in~\cref{sec: AME62 has a common local unitary normal}, we can realize a signature that is locally unitarily equivalent to $f_6$. However, the corresponding local unitary matrices may not be able to realized. This makes us to find the relations between the unitary matrices. Thus, in~\cref{lem:f6-common-orthogonal}, we show that either $\Holant(\mathcal{F})$ is \#P-hard, or a real orthogonal holographic transformation places these matrices projectively in the tetrahedral group. After the transformation, in~\cref{lem:f6-group-alternatives}, we then proved that the projective binary group of the $\Holant$ problem: it is either the tetrahedral group $\Gamma$ or a subgroup of $K_4$. We deal with the two cases in~\cref{lem:f6-tetrahedral-case,lem:f6-klein-case}, both cases give us the realization of $f_6$ and four Bell signatures. 

In \cref{sec: f6 is hard}, we prove hardness when $f_6$ is available.
The proof has two cases.
If every signature in $\mathcal{F}$ has parity, we apply a technique called ``realification" to prove every signature in $\mathcal{F}$ is \emph{real} up to a complex phase, and then apply the real Holant dichotomy (\cref{subsec: f6 parity}).
Otherwise, we reduce the arity while preserving the non-parity property until we obtain a binary signature without parity.
We then extend the ``non-$\mathcal{B}$ hard'' property of $f_6$ in~\cite{realholant} to the ``non-$\Gamma$ hard'' property and complete the proof using a minimal counterexample argument (\cref{subsec: f6 non-parity}).

For $2n=8$, we first realize $f_8$ in \cref{subsec: realize f8 Reed Muller}.
In \cref{subsec: f8 available hardness}, we prove that $\Holant(f_8,\mathcal{F})$ is \#P-hard under the assumption that $\mathcal{F}\not\subseteq\mathscr{A}$ .
The proof follows the framework for real Holant problems~\cite{realholant} and again uses a realification argument similar to that in \cref{subsec: f6 parity}.
Finally, for $2n\geq 10$, 
we can reduce the arity while keeping the resulting signature outside $\mathcal{U}^{\otimes}$, thus completing the induction (\cref{sec: 2n >= 10}).

\section{Second Order Orthogonality}\label{sec: second order orthoganality}
In this section, we will prove that if $\mathcal{F}$ contains an irreducible signature that doesn't satisfy {\sc 2nd-Orth}, then $\Holant(\mathcal{F})$ is \#P-hard (\cref{lm: not 2nd-orth is hard}).
First, \cref{def: k-th order orthogonality} specializes to the following when $k=2$.
\begin{definition}[Second order orthogonality \cite{realholant}]\label{def: second-order-orth}
    Let $f$ be a complex-valued signature of arity $n \geqslant 4$. It satisfies the \emph{second order orthogonality ({\sc 2nd-Orth})} if there exists some $\lambda\neq 0$ such that for all pairs of indices $\{i, j\}\subseteq [n]$, the entries of $f$ satisfy 
    \begin{equation*}
        |{\bf f}_{ij}^{00}|^2=|{\bf f}_{ij}^{01}|^2=|{\bf f}_{ij}^{10}|^2=|{\bf f}_{ij}^{11}|^2=\lambda, 
    \text{ ~~~~\rm and }~~~~ \langle{\bf f}_{ij}^{ab}, {\bf f}_{ij}^{cd}\rangle =0 ~~~~\text{\rm  for  all } (a, b)\neq (c, d).
    \end{equation*}
\end{definition}

Similar to \cref{lm: 1st-Orth uniform}, one can easily prove:
\begin{lemma}\label{lm: 2nd-Orth uniform}
    Let $f$ be a signature of arity $n$.
    If for all indices $i\neq j\in [n]$, $M(\mathfrak{m}_{ij}f)= \lambda_{ij} I_4$ for some real $\lambda_{ij}\neq 0$, then $f$ satisfies {\sc 2nd-Orth} (i.e., all $\lambda_{ij}$ have the same value).
\end{lemma}


We first record a consequence of {\sc 2nd-Orth} that will be frequently used in the subsequent sections.
\begin{lemma}\label{lm: 2nd-orth contraction nonzero}
    Suppose $f$ satisfies {\sc 2nd-Orth}, and $b$ is a non-zero binary signature.
    Then for every pair of indices $\{i,j\}$, $\partial_{ij}^bf$ is a non-zero signature.
    Moreover, $|\partial_{ij}^bf|^2=\lambda|b|^2$, where $\lambda=|f_{ij}^{00}|^2=|f_{ij}^{01}|^2=|f_{ij}^{10}|^2=|f_{ij}^{11}|^2>0.$
\end{lemma}
\begin{proof}
    $$|\partial_{ij}^{b}f|^2
    %=\braket{\sum_{r,s\in\{0,1\}}b(r,s)f_{ij}^{rs},\sum_{r,s\in\{0,1\}}b(p,q)f_{ij}^{pq}}
    =
    \sum_{r,s,p,q\in\{0,1\}}
    b(r,s)\overline{b(p,q)}
    \langle f_{ij}^{rs},f_{ij}^{pq}\rangle=\sum_{r,s\in\{0,1\}}|b(r,s)|^2\braket{f_{ij}^{rs},f_{ij}^{rs}}=|b|^2\cdot\lambda$$
    by {\sc 2nd-Orth}, where $\lambda=|f_{ij}^{00}|^2=|f_{ij}^{01}|^2=|f_{ij}^{10}|^2=|f_{ij}^{11}|^2>0$.
\end{proof}

Recall $\mathfrak{m}_{ij}f$ is the 4-arity signature representing the mating of $f$ and $\overline{f}$ by their corresponding variables except for $x_i$ and $x_j$.
Note that $f$ satisfies {\sc 2nd-Orth} iff $M(\mathfrak{m}_{ij}f)=\lambda I_4$ for some $\lambda>0$.
In the setting of $\holant{\neq_2}{\hF}$, the mating operation corresponds to connecting variables except for $x_i$ and $x_j$ of $\widehat{f}$ and $\widehat{\overline{f}}$ using $\neq_2$.
We denote the resulting signature by $\widehat{\mathfrak m}_{ij} \widehat{f}$.
Clearly $\widehat{\mathfrak m}_{ij} \widehat{f}=\widehat{\mathfrak{m}_{ij}f}$, so $f$ satisfies {\sc 2nd-Orth} iff $M(\widehat{\mathfrak m}_{ij} \widehat{f})=\lambda N_4$ for some $\lambda>0$.

In the following, we assume that $\mathcal{F}$ is conjugate closed and doesn't satisfy condition \rm{(\ref{cond: tractable classes})}.
Recall that $\hF$ is arrow reversal closed by \cref{lm: conjugate closed and ars closed}.
We want to show that every irreducible signature of arity $2n\ge 4$ must satisfy {\sc 2nd-Orth}, otherwise $\Holant(\mathcal{F})$ is \#P-hard.
First we can realize the 4-ary signature $\mathfrak{m}_{ij}f$ for every pair of indices $\{i,j\}$.
We consider all possible cases that $\mathfrak{m}_{ij}f$ falls into the tractable classes of \cref{thm: single quaternary dichotomy}.

\subsection{Unary and Binary Tensor}
We first consider the case that $\mathfrak{m}_{ij}f\in \braket{\mathcal{T}}$ for some pair of indices $\{i,j\}.$

\begin{lemma}\label{lm: 2nd-orth mating reducible}
    Let $f\in \mathcal{F}$ be a signature of arity $2n\ge 4.$
    If there exists pair of indices $\{i,j\}$ such that $\mathfrak{m}_{ij}f\in \braket{\mathcal{T}}$, then one of the following holds:
    \begin{itemize}
        \item $\Holant(\mathcal{F})$ is \#P-hard.
        \item $f$ is reducible.
        \item $M(\mathfrak{m}_{ij}f)=\lambda_{ij}I_4$ for some $\lambda_{ij}>0.$
    \end{itemize}
\end{lemma}
\begin{proof}
    If $\mathfrak{m}_{ij}f$ fails {\sc 1st-Orth}, then $\Holant(\mathcal{F})$ is \#P-hard by \cref{lm: not 1st-orth is hard}.
    Now we assume that $\mathfrak{m}_{ij}f$ satisfies {\sc 1st-Orth}.
    If $\mathfrak{m}_{ij}f$ has a unary factor $u$, $u$ is realizable by \cref{lm: decomposition}.
    By \cref{lm: odd holant with conjugation}, $\Holant(\mathcal{F}
    )$ is \#P-hard.
    In the following, assume that $\mathfrak{m}_{ij}f=b_1\otimes b_2$ for two binary signatures $b_1$ and $b_2$.
    Let $R=M(\mathfrak{m}_{ij}f),\,U=M(b_1)$ and $V=M(b_2)$.
    Since $\mathfrak{m}_{ij}f=b_1\otimes b_2$, there are three possibilities.
    \begin{itemize}
        \item $R_{ab,cd}=U_{ab}V_{cd}$ for every $\{a,b,c,d\}\in\{1,2\}.$
        Then $\mathrm{rank}(R)=1.$
        Since $\mathrm{rank}(R)=\mathrm{rank}(M(\mathfrak{m}_{ij}f))=\mathrm{rank}(M_{ij}(f)M_{ij}(f)^\dagger)=\mathrm{rank}(M_{ij}f)$, we have $\mathrm{rank}(M_{ij}f)=1.$
        Therefore, $f$ factors into two binaries.
        Thus $f$ is reducible.
        
        \item $R_{ab,cd}=U_{ac}V_{bd}$ for every $\{a,b,c,d\}\in\{1,2\}.$
        Then $R=U\otimes V.$
        Recall that $R_{ab,cd}=\braket{\widehat{\mathbf{f}}_{ij}^{ab},\widehat{\mathbf{f}}_{ij}^{cd}}.$
        %By \cref{lm: not 1st-orth is hard}, either $\Holant(\neq_2\mid \widehat{\mathcal{F}})\equiv_T\Holant(\mathcal{F})$ is \#P-hard, or  
        By {\sc 1st-Orth} at $x_i$, there exist some $\mu>0$ such that $\mu \delta_{ac}=\braket{\widehat{\mathbf{f}_i^a},\widehat{\mathbf{f}_i^c}}=R_{a1,c1}+R_{a2,c2}=U_{ac}(V_{11}+V_{22})$, for every $a,c\in\{1,2\}.$
        Set $s=V_{11}+V_{22}$.
        Taking $a=c=1$ gives $\mu=U_{11}s>0$, so $s\neq 0.$
        Therefore $U_{11}=U_{22}=\frac{\mu}{s}$ and $U_{12}=U_{21}=0$.
        Thus $U=uI_2$ for some $u\neq 0.$
        Similarly, {\sc 1st-Orth} at $x_j$ gives $\mu\delta_{bd}=(U_{11}+U_{22})V_{bd}=2uV_{bd}.$
        Hence $V=\frac{\mu}{2u}I_2.$
        So $R=U\otimes V=\frac{\mu}{2}I_4$, $\mu>0.$
        It follows that $M(\widehat{\mathfrak{m}}_{ij}\widehat{f})= \frac{\mu}{2} N_4$.
        Taking $\lambda_{ij}=\frac{\mu}{2}$ completes this case.
    
        \item $R_{ab,cd}=U_{ad}V_{bc}$ for every $\{a,b,c,d\}\in\{1,2\}.$
        {\sc 1st-Orth} at $x_i$ says there exists some $\mu>0$, \(\mu\delta_{ac}=R_{a1,c1}+R_{a2,c2}=U_{a1}V_{1c}+U_{a2}V_{2c}=(UV)_{ac}\), for every $\{a,c\}\in\{1,2\}.$
        Thus $UV=\mu I_2,\,\det(U)\det(V)=\mu^2.$
        Consider $R'=U\otimes V$, so $R'_{ab,cd}=U_{ac}V_{bd}.$
        Exchanging the two middle columns of $R'$ indexed by $12$ and $21$ switches the column index $c$ and $d$.
        So we obtain $U_{ad}V_{bc}=R_{ab,cd}.$
        Thus, $\det(R)=-\det(R')=-\det(U)^2\det(V)^2=-\mu^4<0.$
        However, $\det(R)=\det(M(\mathfrak{m}_{ij}f))=\det(M_{ij}(f)M_{ij}(f)^\dagger)\ge0$, contradiction.
        So this case cannot happen.
    \end{itemize}
    In all above cases, either $f$ is reducible, or $M(\mathfrak{m}_{ij}f)=\lambda_{ij}I_4$ for some $\lambda_{ij}>0.$
    The lemma is proved.
\end{proof}

In the following, we may assume that $\mathfrak{m}_{ij}f$ is irreducible for every pair of indices $\{i,j\}$.

\subsection{Affine, Product, Local-affine Transformable}

In this subsection, we consider the case that some irreducible $\mathfrak{m}_{ij}f$ is $\mathscr{A}$- or $\mathscr{P}$- or $\mathscr{L}$-transformable.
We prove either $\Holant(\mathcal{F})$ is \#P-hard or $M(\mathfrak{m}_{ij}f)=\lambda_{ij}I_4$ for some $\lambda_{ij}>0.$

\begin{lemma}\label{lm: APL-transformable singular value}
    Let $h$ be an irreducible signature of arity 4, and $h$ satisfies {\sc 1st-Orth}.
    If $h$ is $\mathscr{A}$- or $\mathscr{P}$- or $\mathscr{L}$-transformable, then for any $i,j\in[4]$, the 4 by 4 signature matrix $M_{ij}(h)$ has equal non-zero singular values and $\mathrm{rank}(M_{ij}(h))\in \{2,4\}.$
\end{lemma}
\begin{proof}
    By \cref{lm: APL-transformable}, there exists $T_1,T_2,T_3,T_4\in \mathrm{GL}_2(\C)$ such that $h=(T_1\otimes T_2 \otimes T_3 \otimes T_4)a$, where $a\in \mathscr{A}.$
    By \cref{lm: affine is 1st-orth}, $a$ satisfies {\sc 1st-Orth}.
    Consider the singular value decomposition for every $T_i, \,1\le i\le 4$.
    We have $T_i=U_i\mathrm{diag}(\sigma_{i,1},\sigma_{i,2})V_i^\dagger$, where $U_i,V_i$ are 2 by 2 unitary matrices, and $\sigma_{i,1}\ge\sigma_{i,2}>0$ are the singular values of $T_i.$
    Let $c_i=\sqrt{\sigma_{i,1}\cdot \sigma_{i,2}}$, $r_i=\sqrt{\sigma_{i,1}/\sigma_{i,2}}$, $D_i=\mathrm{diag}(r_i,r_i^{-1})$, we have $T_i=c_iU_i D_i V_i^\dagger$ for $1\le i\le 4.$
    Let $c=c_1c_2c_3c_4$. 
    Then
    $$
    h=c(U_1\otimes \cdots \otimes U_4)(D_1\otimes \cdots \otimes D_4)(V_1^\dagger \otimes \cdots \otimes V_4^\dagger)a.
    $$
    Let $f=(V_1^\dagger\otimes \cdots\otimes V_4^\dagger)a$ and $g=c^{-1}(U_1^\dagger\otimes \cdots \otimes U_4^\dagger)h$, we have
    \begin{equation}\label{eq: g = Df}
        g=(D_1\otimes \cdots \otimes D_4)f,\qquad D_i=\mathrm{diag}(r_i,r_i^{-1}).
    \end{equation}
    Since scalar and binary extension by unitary matrices do not change {\sc 1st-Orth}, both $f$ and $g$ satisfies {\sc 1st-Orth}.
    
    We next show $f=g$.
    Let $x=(x_1,x_2,x_3,x_4)\in \{0,1\}^{4}.$
    By \eqref{eq: g = Df}, $g(x)=d(x)f(x)$, where $d(x)=\prod_{i=1}^4r_i^{1-2x_i}>0$.
    Because $f$ satisfies {\sc 1st-Orth}, $||f_i^0||^2=||f_i^1||^2$ for every $1\le i\le 4$.
    So $\sum_{x\in\{0,1\}^4,\,x_i=0}|f(x)|^2=\sum_{x\in\{0,1\}^4,\,x_i=1}|f(x)|^2$, for every $1\le i\le 4.$
    Define $p(x)=|f(x)|^2/||f||^2$ for every $x\in\{0,1\}^4.$
    Then 
    $$
    \sum_{x\in\{0,1\}^4}(1-2x_i)p(x)=\sum_{x\in\{0,1\}^4,x_i=0}p(x)-\sum_{x\in\{0,1\}^4,x_i=1}p(x)=0,\quad \forall i, 1\le i\le 4.
    $$
    Thus
    \begin{equation}\label{eq: ||g|| over ||f||}
        \begin{aligned}
            \frac{||g||^2}{||f||^2}=&\frac{\sum_{x\in\{0,1\}^4}d(x)^2|f(x)|^2}{||f||^2}
            =\sum_{x\in\{0,1\}^4}p(x)d(x)^2\ge \prod_{x\in\{0,1\}^4}(d(x)^2)^{p(x)}\\
            =&\prod_{x\in\{0,1\}^4}\prod_{i=1}^4(r_i^{1-2x_i})^{2p(x)}
            =\prod_{i=1}^4r_i^{2\sum_{x\in\{0,1\}^4}(1-2x_i)p(x)}=1.
        \end{aligned}
    \end{equation}
    The inequality is by AM-GM.
    By \eqref{eq: ||g|| over ||f||}, we know $||g||^2\ge ||f||^2.$

    Conversely, we have 
    \begin{equation}\label{eq: f = Dg}
        f=(D_1^{-1}\otimes \cdots \otimes D_4^{-1})g,\qquad D_i^{-1}=\mathrm{diag}(r_i^{-1},r_i).
    \end{equation}
    Using the fact that $g$ satisfies {\sc 1st-Orth}, the same argument above gives $||f||^2\ge ||g||^2.$
    So $||f||=||g||,$ and the inequality in AM-GM becomes equality.
    So $d(x)^2=c$ for all $x$ with $p(x)>0$.
    Notice that $1=\prod_{x\in\{0,1\}^4}d(x)^{2p(x)}=\prod_{x\in\{0,1\}^4}c^{p(x)}=c^{\sum_{x\in\{0,1\}^4}p(x)}=c$.
    Also, $p(x)>0\iff x\in \mathscr{S}(f)$.
    So $d(x)=1$ for every $x\in \mathscr{S}(f).$
    Therefore, $f(x)=g(x)$.

    Recall $f=(V_1^\dagger\otimes \cdots\otimes V_4^\dagger)a$ and $g=c^{-1}(U_1^\dagger\otimes \cdots \otimes U_4^\dagger)h$.
    So $h$ is a scalar multiple of local unitary transformation of \(a\). 
    This doesn't change the rank and singular values of the signature matrix.
    By \cref{lm: affine singular value}, for any $\{i,j\}\in [4]$, $M_{ij}(a)$ has equal non-zero singular values, and $\mathrm{rank}(M_{ij}(a))\in\{1,2,4\}.$
    So $M_{ij}(h)$ has equal non-zero singular values, and $\mathrm{rank}(M_{ij}(h))\in\{1,2,4\}.$

    It remains to exclude that case where $\mathrm{rank}(M_{ij}(h))=1.$
    If so, $M_{ij}(h)=\mathbf{uu^{\tt T}}$ for some $\mathbf{u}\in \mathbb{C}^4.$
    Then $h$ is a tensor product of two binary signatures, both with truth table $\mathbf{u}.$
    This contradicts $h$ is irreducible.
    The lemma is proved.
\end{proof}

The following lemma proves $M(\mathfrak{m}_{ij}f)=\lambda_{ij}I_4$ when $\mathrm{rank}(M(\mathfrak{m}_{ij}f))=4.$

\begin{lemma}\label{lm: APL-trans rank 4}
    Let $f\in \mathcal{F}$ be a signature of arity $2n\ge 4$.
    If there exists a pair of indices $\{i,j\}$ such that (1) $\mathfrak{m}_{ij}f$ is irreducible and satisfies {\sc 1st-Orth}; (2) $\mathfrak{m}_{ij}f$ is $\mathscr{A}$- or $\mathscr{P}$- or $\mathscr{L}$-transformable; (3) $\mathrm{rank}(M(\mathfrak{m}_{ij}f))=4$, then $M(\mathfrak{m}_{ij}f)=\lambda_{ij}I_4$ for some $\lambda_{ij}>0.$
\end{lemma}
\begin{proof}
    Let $h=\mathfrak{m}_{ij}f$.
    Apply \cref{lm: APL-transformable singular value}, $M_{ij}(h)$ has 4 equal singular values $\lambda_{ij}$.
    Because $M_{ij}(h)=M(\mathfrak{m}_{ij}f)=M_{ij}(f)M_{ij}(f)^\dagger$ is Hermitian and positive semi-definite, its eigenvalues and singular values are the same, and are all equal to $\lambda_{ij}\ge 0$.
    Since $\mathrm{rank}(M_{ij}(h))=4$, $\lambda_{ij}>0.$
    By spectral decomposition of Hermitian matrix, there exists unitary matrix $U$ such that  $M_{ij}(h)=U\mathrm{diag}(\lambda_{ij},\lambda_{ij},\lambda_{ij},\lambda_{ij})U^{\dagger}=\lambda_{ij}UU^\dagger=\lambda_{ij}I_4$.
    The lemma is proved.
\end{proof}

In the following, we will show $\Holant(\mathcal{F})$ is \#P-hard if $\mathrm{rank}(M(\mathfrak{m}_{ij}f))=2$ for some pair of indices $\{i,j\}.$

\begin{definition}
Let $R\in\mathbb C^{4\times4}$ have rows and columns indexed
by pairs of bits. 
Its partial traces are the $2\times2$ matrices
defined by
\[
  (\operatorname{Tr}_1R)_{b,d}
  :=\sum_{a\in\{0,1\}}R_{ab,ad},
  \qquad
  (\operatorname{Tr}_2R)_{a,c}
  :=\sum_{b\in\{0,1\}}R_{ab,cb}.
\]
For an arity 4 signature $g$, we write
$\operatorname{Tr}_k(g):=\operatorname{Tr}_k(M(g))$ for $k=1,2$,
where $M(g)_{ab,cd}=g(a,b,c,d)$.
\end{definition}

\begin{lemma}\label{lem:mating-partial-traces}
Let $f$ be a signature of arity $2n\ge 4$ satisfying {\sc 1st-Orth}.
Then there exists some $\mu>0$, for every pair of distinct indices $\{i,j\}$,
\(
  \operatorname{Tr}_1(\mathfrak m_{ij}f)
  =\operatorname{Tr}_2(\mathfrak m_{ij}f)
  =\mu I_2.
\)
\end{lemma}

\begin{proof}
By permuting variables, assume that $(i,j)=(1,2)$.
Since $f$ satisfies {\sc 1st-Orth}, there exists $\mu>0$ such that
$M_1(f)M_1(f)^\dagger=M_2(f)M_2(f)^\dagger=\mu I_2$.

Put $R=M(\mathfrak m_{12}f)$. By the definition of mating,
$R_{ab,cd}=\sum_{x\in\{0,1\}^{2n-2}}
f(a,b,x)\overline{f(c,d,x)}$.
Consequently, $(\operatorname{Tr}_2R)_{a,c}=\sum_{b\in\{0,1\},x\in\{0,1\}^{2n-2}}
f(a,b,x)\overline{f(c,b,x)}$,
which is precisely the $(a,c)$ entry of $M_1(f)M_1(f)^\dagger$.
Thus $\operatorname{Tr}_2R=\mu I_2$.
Summing over the first coordinate instead gives
$\operatorname{Tr}_1R=M_2(f)M_2(f)^\dagger=\mu I_2$,
proving the claim.
\end{proof}

\begin{lemma}\label{lm:rank-two-matrix-form}
Let $h$ be an arity 4 signature.
Suppose that $M(h)$ is Hermitian, positive semi-definite of rank two with two nonzero equal eigenvalues, and
$\operatorname{Tr}_1R=\operatorname{Tr}_2R=\mu I_2$
for some $\mu>0$.
Then there exist traceless Hermitian matrices
$A,B\in\mathbb C^{2\times2}$ with $A^2=B^2=I_2$ such that
\(
  M(h)=\frac{\mu}{2}(I_4+A\otimes B).
\)
\end{lemma}

\begin{proof}
Let $R=M(h)$, and $\lambda>0$ be the common nonzero eigenvalue of $R$.
Since $\mathrm{rank}(R)=2$, we have $\operatorname{Tr}R=2\lambda$.
Also,
$\operatorname{Tr}R=\operatorname{Tr}(\operatorname{Tr}_1R)=2\mu$,
so $\lambda=\mu$.
By the spectral decomposition of Hermitian matrices, we have $R=U\mathrm{diag}(\lambda,\lambda,0,0)U^{\dagger}$ for some unitary matrix $U$.
Hence $R^2=\lambda R=\mu R$.
Let $W:=2R/\mu-I_4$.
It is Hermitian, satisfies $W^2=I_4$,
and has both partial traces zero.

Writing $W$ in $2\times2$ blocks, Hermitian symmetry and
$\operatorname{Tr}_1W=0$ give
\[
  W=\begin{pmatrix}
      D&C\\
      C^\dagger&-D
    \end{pmatrix},
  \qquad D=D^\dagger.
\]
The condition $\operatorname{Tr}_2W=0$ further gives
$\operatorname{Tr}D=\operatorname{Tr}C=0$.
Comparing blocks in $W^2=I_4$, we obtain
$D^2+CC^\dagger=D^2+C^\dagger C=I_2$ and $DC=CD$.
Therefore $C$ is normal and commutes with the Hermitian matrix $D$.

By simultaneous unitary diagonalization, there is a unitary matrix
$V$ such that
$V^\dagger DV=\operatorname{diag}(d,-d)$ and
$V^\dagger CV=\operatorname{diag}(z,-z)$,
where $d\in\mathbb R$ and $z\in\mathbb C$.
Here the opposite diagonal entries follow from the zero traces.
The identity $D^2+CC^\dagger=I_2$ gives $d^2+|z|^2=1$.

Define
\[
  A=\begin{pmatrix}d&z\\\overline z&-d\end{pmatrix},
  \qquad
  B=V\begin{pmatrix}1&0\\0&-1\end{pmatrix}V^\dagger.
\]
Both matrices are traceless and Hermitian.
Moreover, $B^2=I_2$ and
$A^2=(d^2+|z|^2)I_2=I_2$.
Since $D=dB$ and $C=zB$, the block expression for $W$
gives $W=A\otimes B$.
Substituting the definition of $W$ proves the claimed formula.
\end{proof}

\begin{lemma}\label{lm:realize-I-plus-AA}
Let $\mathcal F$ be a set of signatures. 
Suppose an arity 4 signature $h$ is realizable from $\mathcal F$ and
$M(h)=\lambda(I_4+A\otimes B)$, where $\lambda>0$
and $A,B\in\mathbb C^{2\times2}$ are traceless Hermitian matrices
satisfying $A^2=B^2=I_2$.
Then we can realize a signature $g$ with matrix
$M(g)=I_4+A\otimes A$.
\end{lemma}

\begin{proof}
Connect variables $x_2,x_4$ of the first copy of $h$ to variables
$x_4,x_2$ of the second, respectively.
The resulting signature is
$g'(a,b,c,d)=\sum_{x,y\in\{0,1\}}h(a,x,c,y)h(b,y,d,x)$.
Since $\operatorname{Tr}B=0$ and
$\operatorname{Tr}(B^2)=2$, direct expansion gives
\[
\begin{aligned}
g'(a,b,c,d)
&=\lambda^2\sum_{x,y\in\{0,1\}}
  (\delta_{ac}\delta_{xy}+A_{ac}B_{xy})
  (\delta_{bd}\delta_{yx}+A_{bd}B_{yx})\\
&=2\lambda^2
  (\delta_{ac}\delta_{bd}+A_{ac}A_{bd}).
\end{aligned}
\]
Thus $M(g')=2\lambda^2(I_4+A\otimes A)$.
Normalizing by $2\lambda^2$ proves the lemma.
\end{proof}

\begin{lemma}\label{lm: I+AA implies hard}
    Let $\mathcal{F}$ be a conjugate closed set of signatures.
    If from $\mathcal{F}$ we can realize a 4-ary signature $g$ with matrix $M(g)=I_4+A\otimes A$, with $A=A^{\dagger}, A^2=I_2$ and $\mathrm{Tr}(A)=0$, then $\Holant(\mathcal{F})$ is \#P-hard.
\end{lemma}
\begin{proof}
    By assumption, we may write the spectral decomposition of $A$ as $A=U\mathrm{diag}(1,-1)U^\dagger$, where $U=[\mathbf{u}_1, \mathbf{u}_2]$ is a 2 by 2 unitary matrix.
    Consider $U^{\tt T}U=\left[\begin{matrix}
        \mathbf{u}_1^{\tt T}\mathbf{u}_1 & \mathbf{u}_1^{\tt T}\mathbf{u}_2\\
        \mathbf{u}_2^{\tt T}\mathbf{u}_1 & \mathbf{u}_2^{\tt T}\mathbf{u}_2
    \end{matrix}\right]=\left[\begin{matrix}
        q_1 & z\\
        z & q_2
    \end{matrix}\right],$
    where $q_1=\mathbf{u}_1^{\tt T}\mathbf{u}_1,q_2=\mathbf{u}_2^{\tt T}\mathbf{u}_2$ and $z=\mathbf{u}_1^{\tt T}\mathbf{u}_2=\mathbf{u}_2^{\tt T}\mathbf{u}_1.$
    Since $U$ is unitary, $U^{\tt T}U$ is also unitary.
    So $|q_1|^2+|z|^2=|q_2|^2+|z|^2$, $|q_1|=|q_2|.$
    We may choose a complex phase of the eigenvectors $\mathbf{u}_1$ and $\mathbf{u}_2$, such that $q_1=\mathbf{u}_1^{\tt T}\mathbf{u}_1=\mathbf{u}_2^{\tt T}\mathbf{u}_2=q_2=r\ge 0$.
    \begin{itemize}
        \item If $r=0$, then $\mathbf{u}_1^{\tt T}\mathbf{u}_1=\mathbf{u}_2^{\tt T}\mathbf{u}_2=0$.
        So $\{\mathbf{u_1},\mathbf{u_2}\}=\{\frac{1}{\sqrt{2}}[1, \mathfrak{i}]^{\tt T},\frac{1}{\sqrt{2}}[1,-\mathfrak{i}]^{\tt T}\}$.
        So $A=U\mathrm{diag}(1,-1)U^\dagger=\pm Y$, where $Y=\left[\begin{matrix}
            0 & -\mathfrak{i}\\
            \mathfrak{i} & 0
        \end{matrix}\right]$.
        So $M(g)=I_4+Y\otimes Y.$
        Direct computation shows
        \[
        (K^{-1})^{\otimes 2}M(g)((K^{-1})^{\tt T})^{\otimes 2}=\begin{bmatrix}
            0& 0 & 0 & 2\\
            0& 0 & 0 & 0\\
            0& 0 & 0 & 0\\
            2& 0 & 0 & 0\\
        \end{bmatrix}.
        \]
        So $K^{-1}g=2(\neq_4).$
        By a $K$-holographic transformation, we have $\Holant(\neq_2\mid \widehat{\mathcal{F}},\neq_4)\le_T \Holant(\mathcal{F},g)\le_T\Holant(\mathcal{F}).$
        By \cref{lem: In conjugate closed setting disequality 4 gives hardness}, $\Holant(\neq_2\mid \widehat{\mathcal{F}},\neq_4)$ is \#P-hard.
        It follows that $\Holant(\mathcal{F})$ is \#P-hard.
        
        \item Suppose $r>0.$
        Consider merging variables $x_1$ and $x_2$ of $g$ and obtain the signature $b=\partial_{12}g.$
        Since $g(x_1,x_2,x_3,x_4)=\delta_{x_1,x_3}\delta_{x_2,x_4}+A_{x_1,x_3}A_{x_2,x_4}$, we have 
        $b(x_3,x_4)=\partial_{12}g(x_3,x_4)=\sum_{t\in\{0,1\}}g(t,t,x_3,x_4)=\sum_{t\in\{0,1\}}(\delta_{t,x_3}\delta_{t,x_4}+A_{t,x_3}A_{t,x_4})=\delta_{x_3,x_4}+(A^{\tt T}A)_{x_3,x_4}.$
        So $M(b)=I_2+A^{\tt T}A.$
        Let $B=M(b).$
        We have
        \begin{equation}
            \begin{aligned}
                U^{\tt T}BU
                =&U^{\tt T}U+U^{\tt T}A^{\tt T}AU\\
                =&U^{\tt T}U+U^{\tt T}\overline{U}\mathrm{diag}(1,-1)U^{\tt T}U\mathrm{diag}(1,-1)U^{\dagger}\\
                =& \left[\begin{matrix}
                r & z\\
                z & r
                \end{matrix}\right]
                +\mathrm{diag}(1,-1)\left[\begin{matrix}
                r & z\\
                z & r
                \end{matrix}\right]\mathrm{diag}(1,-1)\\
                =&\left[\begin{matrix}
                2r & 0\\
                0 & 2r
                \end{matrix}\right].
            \end{aligned}
        \end{equation}
        Thus, $B=2r\overline{U}U^\dagger.$
        Let $c=\frac{b}{2r}$ and $C=M(c)$.
        Since $\mathcal{F}$ is conjugate closed, $\overline{c}$ is also realizable, with matrix $\overline{C}=UU^{\tt T}.$
        %So $\overline{c}=U^{\otimes 2}(=_2).$
        Connecting $c$ and $\overline{c}$ using $=_2$, we can realize the signature with matrix $C\overline{C}=\overline{U} U^\dagger UU^{\tt T}=I_2$, i.e., we can realize $=_2$ on RHS.
        Now connect two copies of $\overline{c}$ with variables $x_3$ and $x_4$ of $g$ respectively.
        This will give the 4-ary signature $e=2\cdot U^{\otimes}(=_4)$.
        To see this, first notice that $M(g)=I_2\otimes I_2+A\otimes A=(\mathbf{u}_1\mathbf{u}_1^\dagger+\mathbf{u}_2\mathbf{u}_2^\dagger)^{\otimes 2}+(\mathbf{u}_1\mathbf{u}_1^\dagger-\mathbf{u}_2\mathbf{u}_2^\dagger)^{\otimes 2}=2\sum_{s=0}^1\mathbf{u}_s\otimes \mathbf{u}_s\otimes \overline{\mathbf{u}_s}\otimes \overline{\mathbf{u}_s}$.
        Then $e=(I_2\otimes I_2\otimes \overline{C}\otimes\overline{C})g=2(\mathbf{u}_0^{\otimes 4}+\mathbf{u}_1^{\otimes 4})=2\cdot U^{\otimes}(=_4).$
        On the other hand, $c$ is realizable on LHS by connecting two copies of $=_2$ on LHS with one copy of $c$ on RHS.
        Therefore,
        \[
        \Holant(c\mid \mathcal{F}\cup\{=_2,e\})\le_T\Holant(\mathcal{F}).
        \]
        Consider the holographic transformation by $U$.
        We have
        \[
        \Holant(=_2\mid U^\dagger (\mathcal{F}\cup\{=_2\}),=_4)
        \le_T
        \Holant(c\mid \mathcal{F}\cup\{=_2,e\})\le_T\Holant(\mathcal{F}).
        \]
        By \cref{lm: CSP2 to holant =4}, $\mathrm{CSP}_2({U^\dagger}(\mathcal{F}\cup\{=_2\}))\le_T\Holant(=_2\mid U^\dagger (\mathcal{F}\cup\{=_2\}),=_4).$
        By \cref{lm: CSP with =2}, $\mathrm{CSP}_2({U^\dagger}(\mathcal{F}\cup\{=_2\}))$ is \#P-hard.
        It follows that $\Holant(\mathcal{F})$ is \#P-hard.
    \end{itemize}
    
\end{proof}

\begin{lemma}\label{lm: APL-trans rank 2}
    Let $f\in \mathcal{F}$ be a signature of arity $2n\ge 4$.
    If there exists a pair of indices $\{i,j\}$ such that (1) $\mathfrak{m}_{ij}f$ is $\mathscr{A}$- or $\mathscr{P}$- or $\mathscr{L}$-transformable; (2) $\mathfrak{m}_{ij}f$ is irreducible and satisfies {\sc 1st-Orth}; (3) $\mathrm{rank}(M(\mathfrak{m}_{ij}f))=2$, then $\Holant(\mathcal{F})$ is \#P-hard.
\end{lemma}
\begin{proof}
    If $f$ doesn't satisfy {\sc 1st-Orth}, then $\Holant(\mathcal{F})$ is \#P-hard by \cref{lm: not 1st-orth is hard}.
    Now assume that $f$ satisfies {\sc 1st-Orth}.
    By \cref{lem:mating-partial-traces}, there exists $\mu>0$ such that $\mathrm{Tr}_1(\mathfrak{m}_{ij}f)=\mathrm{Tr}_2(\mathfrak{m}_{ij}f)=\mu I_2.$
    
    We realize $h:=\mathfrak{m}_{ij}f$ by mating.
    Since $h$ is irreducible, satisfies {\sc 1st-Orth}, and is $\mathscr{A}$- or $\mathscr{P}$- or $\mathscr{L}$-transformable, 
    \cref{lm: APL-transformable singular value} applies.
    Also, $\mathrm{rank}(M(h))=2$ by assumption, so $M(h)$ has two equal non-zero singular values.
    Notice $M(h)=M_{ij}(f)M_{ij}(f)^\dagger$ is Hermitian and positive semi-definite, so it has two equal non-zero eigenvalues.
    Then \cref{lm:rank-two-matrix-form} applies.
    We have $M(h)=\frac{\mu}{2}(I_4+A\otimes B)$, where $A,B\in \C^{2\times 2}$ are traceless Hermitian matrices with $A^2=B^2=I_2.$
    By \cref{lm:realize-I-plus-AA}, we can realize an arity 4 signature $g$ with matrix $M(g)=I_4+A\otimes A.$
    By \cref{lm: I+AA implies hard}, $\Holant(\mathcal{F})$ is \#P-hard.
    The lemma is proved.
\end{proof}

\begin{lemma}\label{lm: mating is APL-trans implies 2nd-orth}
    Let $f\in \mathcal{\mathcal{F}}$ be a signature of arity $2n\ge 4$.
    If there exists a pair of indices $\{i,j\}$ such that $\mathfrak{m}_{ij}f$ is $\mathscr{A}$- or $\mathscr{P}$- or $\mathscr{L}$-transformable, then either $f$ is reducible, or $\Holant(\mathcal{F})$ is \#P-hard, or $M(\mathfrak{m}_{ij}f)=\lambda_{ij}I_4$ for some $\lambda_{ij}>0.$
\end{lemma}
\begin{proof}
    If $\mathfrak{m}_{ij}f$ has a unary factor, $\Holant(\mathcal{F})$ is \#P-hard by \cref{lm: odd holant with conjugation}.
    If $\mathfrak{m}_{ij}f$ is a tensor product or two binary signatures, the lemma is proved by \cref{lm: 2nd-orth mating reducible}.
    Next we assume that $\mathfrak{m}_{ij}f$ is irreducible.
    
    If $\mathfrak{m}_{ij}f$ doesn't satisfy {\sc 1st-Orth}, then $\Holant(\mathcal{F})$ is \#P-hard by \cref{lm: not 1st-orth is hard}.
    Now we assume that $\mathfrak{m}_{ij}f$ satisfies {\sc 1st-Orth}.
    By \cref{lm: APL-transformable singular value}, $\mathrm{rank}(M(\mathfrak{m}_{ij}f))=2$ or $4$.
    If $\mathrm{rank}(M(\mathfrak{m}_{ij}f))=2$, then $\Holant(\mathcal{F})$ is \#P-hard by \cref{lm: APL-trans rank 2}.
    If $\mathrm{rank}(M(\mathfrak{m}_{ij}f))=4$, then $M(\mathfrak{m}_{ij}f)=\lambda_{ij}I_4$ for some $\lambda_{ij}>0$ by \cref{lm: APL-trans rank 4}.
\end{proof}

\subsection{Weighted Matching and $\mathscr{H}$ Type}


We next show that the case $\mathfrak{m}_{ij}f\in \braket{K\mathcal{M}}$ or $\mathfrak{m}_{ij}f\in \braket{KX\mathcal{M}}$ cannot happen.

\begin{lemma}\label{lm: mating is self inversive}
    Let $h=\mathfrak{m}_{ij}f$.
    Then $\widehat{h}(a,b,c,d)=\overline{\widehat{h}(\overline{c},\overline{d},\overline{a},\overline{b})}$, for every $a,b,c,d\in \{0,1\}.$
\end{lemma}
\begin{proof}
    
    Recall in the mating gadget $\widehat{\mathfrak{m}}_{ij}\widehat{f}$, we connect corresponding variables of $\widehat{f}$ and $\widehat{\overline{f}}$ using $\neq_2.$
    Assign $a,b$ to the two dangling edges on LHS, and $c,d$ to the two dangling edges on RHS.
    Now reflect the gadget by connecting $\widehat{\overline{f}}$ and $\widehat{f}$ using $\neq_2$, and flip the inputs on all dangling edges.
    The inputs change from $(a,b,c,d)$ to $(\overline{c},\overline{d},\overline{a},\overline{b}).$
    Since $\widehat{\overline{f}}(x)=\overline{\widehat{f}(\overline{x})}$ by \cref{lm: conjugation and K-trans},
    in the reversed gadget, the evaluation of $\widehat{\overline{f}}(x)$ on LHS equals to the conjugation of $\widehat{f}(\overline{x})$, while the evaluation of $\widehat{f}(x)$ on RHS equals to the conjugation of $\widehat{\overline{f}}(x)$.
    That is, $\widehat{h}(a,b,c,d)=\overline{\widehat{h}(\overline{c},\overline{d},\overline{a},\overline{b})}$, for every $a,b,c,d\in \{0,1\}.$
\end{proof}

\begin{lemma}\label{lm: 2nd-Orth M and XM cannot happen}
    Let $f\in \mathcal{F}$ be a signature of arity $2n\ge 4$.
    For an arbitrary pair of indices $\{i,j\}$, if $\mathfrak{m}_{ij}f$ is irreducible, 
    then $\mathfrak{m}_{ij}f\notin \braket{K\mathcal{M}}$ and $\mathfrak{m}_{ij}f\notin \braket{KX\mathcal{M}}$.
\end{lemma}


\begin{proof}
    Suppose for a contradiction that $\mathfrak{m}_{ij}f\in \braket{K\mathcal{M}}$ or $\mathfrak{m}_{ij}f\in \braket{KX\mathcal{M}}$.
    Let $h=\mathfrak{m}_{ij}f$.
    Then $\widehat{h}\in \braket{\mathcal{M}}$ or $\widehat{h}\in \braket{X\mathcal{M}}.$
    Since $h$ is irreducible, $\widehat{h}$ is also irreducible.
    So $\widehat{h}\in \mathcal{M}$ or $\widehat{h}\in X\mathcal{M}.$
    Thus $\widehat{h}$ is supported on Hamming weight at most 1 or at least 3 by the definition of $\mathcal{M}$ and $X\mathcal{M}$.
    
    
    If $\mathscr{S}(\widehat{h})$ contains an element of Hamming weight 1 (resp. 3), then by \cref{lm: mating is self inversive}, $\mathscr{S}(\widehat{h})$ also contains an element of Hamming weight 3 (resp. 1).
    This contradicts $\widehat{h}\in {\mathcal{M}}$ or $\widehat{h}\in {X\mathcal{M}}.$
    Therefore, $\widehat{h}=0$.
    This contradicts $\widehat{h}$ is irreducible.
\end{proof}


Finally, we deal with the case that some $\mathfrak{m}_{ij}f\in \mathscr{H}.$

\begin{lemma}\label{lm: 2nd-orth H type}
    Let $f\in \mathcal{F}$ be a signature of arity $2n\ge 4$.
    For an arbitrary pair of indices $\{i,j\}$, if $\mathfrak{m}_{ij}f\in \mathscr{H}$, then either $f$ is reducible, or $\Holant(\mathcal{F})$ is \#P-hard, or $M(\mathfrak{m}_{ij}f)=\lambda_{ij}f$ for some $\lambda_{ij}>0.$
\end{lemma}
\begin{proof}
    Let $h=\mathfrak{m}_{ij}f$ and $\widehat{h}=\widehat{\mathfrak{m}}_{ij}\widehat{f}.$
    By the definition of $\mathscr{H}$, $\widehat{h}$ is supported on $\mathrm{HW}^{\ge}$ or $\mathrm{HW}^{\le}$.
    By \cref{lm: mating is self inversive}, $\widehat{h}(a,b,c,d)=\overline{\widehat{h}(\overline{c},\overline{d},\overline{a},\overline{b})}$, for every $a,b,c,d\in \{0,1\}.$
    So $\widehat{h}$ is supported on Hamming weight $2.$
    Therefore, $\Holant(\neq_2\mid \widehat{h})$ is a six-vertex model.
    By \cref{thm: six vertex model}, either
    $\Holant(\neq_2\mid \widehat{h})\le_T\Holant(\mathcal{F})$ is \#P-hard, or $\widehat{h}\in \mathscr{A}$, or $\widehat{h}\in \mathscr{P}$, or $\widehat{h}\in \braket{\mathcal{M}}$, or $\widehat{h}\in \braket{X\mathcal{M}}$.
    If $\widehat{h}\in \mathscr{A}$ or $\widehat{h}\in \mathscr{P}$,
    then $h$ is $\mathscr{A}$- or $\mathscr{P}$-transformable, and the lemma follows from \cref{lm: mating is APL-trans implies 2nd-orth}.
    Also, the case $\widehat{h}\in \braket{\mathcal{M}}$ or $\widehat{h}\in \braket{X\mathcal{M}}$ cannot happen by \cref{lm: 2nd-Orth M and XM cannot happen}.
    The lemma is proved.
\end{proof}

\subsection{Proof of Second Order Orthogonality}

\iffalse
\begin{lemma}\label{lm: pre not 2nd-orth is hard}
    Let $\widehat f\in\widehat{\mathcal{F}}$ be a signature of arity $2n\geqslant 4$. 
    Then, 
    \begin{itemize}

        \item $\holant{\neq_2}{\widehat{\mathcal{F}}}$ is \#P-hard, or

        \item $\widehat{f}$ is reducible, or
        \item  $M(\widehat{\mathfrak{m}}_{ij}\widehat{f})= \lambda_{ij} N_4$ for some real $\lambda_{ij}\neq 0$.
    \end{itemize}
\end{lemma}

\begin{proof}
    If $\widehat{f}\equiv 0$, then the lemma holds trivially since for all $\{i, j\}$ and any $\widehat{b}_{ij}\neq 0$, $\widehat{b}_{ij}(x_i, x_j)\mid \widehat{f}$.
    Thus, we may assume that $f\not\equiv 0$.
  
    If $\widehat{f}$ does not satisfy {\sc 1st-Orth}, then $f$ does not satisfy it by \cref{prop: invariant 1st-Orth}. 
    By \cref{lm: not 1st-orth is hard}, $\holant{\neq_2}{\widehat{\mathcal{F}}}\equiv_T\holant{=_2}{\mathcal{F}}$ is \#P-hard. 
    Thus, we may assume that $\widehat{f}$ satisfies {\sc 1st-Orth}. 
    Then there exists some $\mu>0$, such that for all indices $i$, 
    $$M(\widehat{\mathfrak{m}}_i\widehat{f})
    =\left[\begin{matrix}
    \langle \widehat{{\bf f}}_i^0, \widehat{{\bf f}}_i^1 \rangle & |\widehat{{\bf f}}_i^0|^2\\
    |\widehat{{\bf f}}_i^1|^2 & \langle \widehat{{\bf f}}_i^1, \widehat{{\bf f}}_i^0 \rangle
    \end{matrix}\right]=
    \mu\left[\begin{matrix}
    0 & 1 \\
     1 & 0
    \end{matrix}\right].
    $$
    
    For any variable $x_i$, we may take another arbitrary variable $x_j$ $(j\neq i)$ and partition the sum in the inner product $\langle \widehat{{\bf f}}_i^0, \widehat{{\bf f}}_i^1\rangle=0$ into two sums depending on whether $x_j=0$ or $1$.
    In the following, we denote $\langle \widehat{{\bf f}}_{ij}^{ab}, \widehat{{\bf f}}_{ij}^{cd} \rangle$ 
    by $\braket{ab\mid cd}$.
    We have 
    $$\langle \widehat{{\bf f}}_i^0, \widehat{{\bf f}}_i^1 \rangle
    = 
    \langle \widehat{{\bf f}}_{ij}^{00}, \widehat{{\bf f}}_{ij}^{10} \rangle+\langle \widehat{{\bf f}}_{ij}^{01}, \widehat{{\bf f}}_{ij}^{11}\rangle 
    =0.$$
    Then $\textcolor{black}{\braket{00\mid 10}}+\textcolor{black}{\braket{01\mid 11}}=0$.
    Note that by exchanging $i$ and $j$ we also have $\textcolor{black}{\braket{00\mid 01}}+\textcolor{black}{\braket{10\mid 11}}=0.$
    Similarly, $\braket{\widehat{\bf f}_i^1, \widehat{\bf f}_i^0}=0$ implies that $\textcolor{black}{\braket{10\mid 00}}+\textcolor{black}{\braket{11\mid 01}}=0$ and $\textcolor{black}{\braket{01\mid 00}}+\textcolor{black}{\braket{11\mid 10}}=0$.
    By $|\widehat{{\bf f}}_i^0|^2=|\widehat{{\bf f}}_i^1|^2,$ we have $\braket{00\mid 00}+\braket{01\mid 01}=\braket{10\mid 10}+\braket{11\mid 11}.$
    Similarly by $|\widehat{{\bf f}}_j^0|^2=|\widehat{{\bf f}}_j^1|^2,$ we have $\braket{00\mid 00}+\braket{10\mid 10}=\braket{01\mid 01}+\braket{11\mid 11}.$
    So $\braket{00\mid 00}=\braket{11\mid 11}=c$ and $\braket{01\mid 01}=\braket{10\mid 10}=d$, where $c,d\in \mathbb{R}_{\ge 0}.$ 
    
    Now, consider $\widehat{\mathfrak{m}}_{ij}\widehat{f}$. 
    We have

$$M(\widehat{\mathfrak m}_{ij}\widehat f)=\begin{bmatrix}
{\widehat{{\bf f}}}_{ij}^{00}\\
{\widehat{{\bf f}}}_{ij}^{01}\\
{\widehat{{\bf f}}}_{ij}^{10}\\
{\widehat{{\bf f}}}_{ij}^{11}\\
\end{bmatrix}
N_2^{\otimes{(2n-2)}}
\left[\begin{matrix}
\widehat{\overline{{\bf f}_{ij}^{00}}}^{\tt T}&
\widehat{\overline{{\bf f}_{ij}^{01}}}^{\tt T}&
\widehat{\overline{{\bf f}_{ij}^{10}}}^{\tt T}&
\widehat{\overline{{\bf f}_{ij}^{11}}}^{\tt T}
\end{matrix}\right].$$
By \cref{lm: conjugation and K-trans},
we have
\begin{equation*}
    N_2^{\otimes (2n-2)}
    \begin{bmatrix}
        \widehat{\overline{{\bf f}_{ij}^{00}}}^{\tt T} &
        \widehat{\overline{{\bf f}_{ij}^{01}}}^{\tt T} &
        \widehat{\overline{{\bf f}_{ij}^{10}}}^{\tt T} &
        \widehat{\overline{{\bf f}_{ij}^{11}}}^{\tt T} 
    \end{bmatrix}
    =
    \begin{bmatrix}
        {\overline{\widehat{{\bf f}}_{ij}^{11}}}^{\tt T}&
        {\overline{\widehat{{\bf f}}_{ij}^{10}}}^{\tt T}&
        {\overline{\widehat{{\bf f}}_{ij}^{01}}}^{\tt T}&
        {\overline{\widehat{{\bf f}}_{ij}^{00}}}^{\tt T}
    \end{bmatrix}.
\end{equation*}
    Then
    $$M(\widehat{\mathfrak m}_{ij}\widehat f)=\begin{bmatrix}
    {\widehat{{\bf f}}}_{ij}^{00}\\
    {\widehat{{\bf f}}}_{ij}^{01}\\
    {\widehat{{\bf f}}}_{ij}^{10}\\
    {\widehat{{\bf f}}}_{ij}^{11}\\
    \end{bmatrix}
    \left[\begin{matrix}
    {\overline{{\widehat{{\bf f}}}_{ij}^{11}}}^{\tt T} &{\overline{{\widehat{{\bf f}}}_{ij}^{10}}}^{\tt T}
    &{\overline{{\widehat{{\bf f}}}_{ij}^{01}}}^{\tt T}&{\overline{{\widehat{{\bf f}}}_{ij}^{00}}}^{\tt T}
    \end{matrix}\right]
    =
    \left[\begin{matrix}
    \braket{00\mid 11} & \textcolor{black}{\braket{00\mid 10}} & 
    \textcolor{black}{\braket{00\mid 01}} & \braket{00\mid 00}\\
    
    \textcolor{black}{\braket{01\mid 11}} & \braket{01\mid 10} & 
    \braket{01\mid 01} & \textcolor{black}{\braket{01\mid 00}}\\

    \textcolor{black}{\braket{10\mid 11}} & \braket{10\mid 10} & 
    \braket{10\mid 01} & \textcolor{black}{\braket{10\mid 00}}\\

    \braket{11\mid 11} & \textcolor{black}{\braket{11\mid 10}} & 
    \textcolor{black}{\braket{11\mid 01}} & \braket{11\mid 00}
    \end{matrix}\right].$$
    
    By \cref{thm: single quaternary dichotomy},
    if $\widehat{\mathfrak m}_{ij}\widehat f=\widehat{\mathfrak{m}_{ij}f}$ belongs to the additional tractable case beyond that of the eight-vertex model, then $\mathfrak{m}_{ij}f\in\braket{K\mathcal{M}}$ or $\mathfrak{m}_{ij}f\in\braket{KX\mathcal{M}}$ or $\mathfrak{m}_{ij}f\in \mathscr{H}$ or $\mathfrak{m}_{ij}f\in \langle\mathscr{T}\rangle.$
    
    
    In the following we consider the first three cases $\mathfrak{m}_{ij}f\in\braket{K\mathcal{M}}$ or $\mathfrak{m}_{ij}f\in\braket{KX\mathcal{M}}$ or $\mathfrak{m}_{ij}f\in \mathscr{H}$.
    In the first two cases, $\widehat{\mathfrak{m}}_{ij}\widehat{f}\in \braket{\mathcal{M}}$ or $\widehat{\mathfrak{m}}_{ij}\widehat{f}\in \braket{X\mathcal{M}}$.
    the entries of weight (the number of ones in $\braket{ab\mid cd}$) 1 (or 3) in the above matrix $M(\widehat{\mathfrak m}_{ij}\widehat f)$ are all 0.
    Combined with the equalities obtained above, $\widehat{\mathfrak{m}}_{ij}\widehat{f}$ has even parity, and it represents a signature of the eight-vertex model.
    Note that $|\braket{ab\mid cd}|^2\le \braket{ab
    \mid ab}\cdot \braket{cd\mid cd}$ by Cauchy-Schwarz inequality.

    If there exists a pair of indices $\{i, j\}$ such that $\Holant(\neq_2 \mid \widehat{\mathfrak{m}}_{ij}\widehat f)$ is \#P-hard, then we are done since $\Holant(\neq_2 \mid \widehat{\mathfrak{m}}_{ij}\widehat f)\leqslant_T \Holant(\neq_2 \mid \widehat{\mathcal{F}})$.
    Thus, we may assume all $\widehat{\mathfrak{m}}_{ij}\widehat f$ belong to the tractable family for eight-vertex models.
    Clearly, by observing its antidiagonal entries of the matrix $M(\widehat{\mathfrak{m}}_{ij}\widehat{f})$, we have $\widehat{\mathfrak{m}}_{ij}\widehat{f} \not\equiv 0$ since $\widehat f\not \equiv 0$.
    By \cref{thm: eight-vertex model},  there are three possible cases.

\begin{itemize}
    \item  $\widehat{\mathfrak{m}}_{ij}\widehat f$ has support of size $2$. 
    Since $f\not\equiv0$, the anti-diagonals of $\widehat{\mathfrak{m}}_{ij} \widehat{f}$ contain at least one non-zero entry.
    According to the positions of non-zero entries in $M(\widehat{\mathfrak{m}}_{ij}\widehat f)$, there are two possible sub-cases:
    (1) $\widehat{\mathfrak{m}}_{ij}\widehat f$ can be factorized as the tensor product of a unary signature and a ternary signature.
    Then by \cref{lm: decomposition} and \cref{lm: odd holant with conjugation}, $\holant{\neq_2}{\widehat{\mathfrak{m}}_{ij}\widehat f}\le_T\Holant(\neq_2 \mid \widehat{\mathcal{F}})$ is \#P-hard.
    (2) After a possible permutation of variables, we can realize a weighted $\neq_4$: $ \left[\begin{smallmatrix}
    0 & 0 & 0 & a\\
     0 & 0 & 0 & 0\\
      0 & 0 & 0 & 0\\
       b & 0 & 0 & 0\\
    \end{smallmatrix}\right]$ where $a,b> 0.$
    Similar to \cref{lm: 2 by 2 interpolation}, we can do polynomial interpolation and realize $(\neq_4)=\left[\begin{smallmatrix}
    0 & 0 & 0 & 1\\
     0 & 0 & 0 & 0\\
      0 & 0 & 0 & 0\\
       1 & 0 & 0 & 0\\
    \end{smallmatrix}\right]$.
    By~\cref{lem: In conjugate closed setting disequality 4 gives hardness}, $\holant{\neq_2}{\hF}$ is $\#P$-hard.
    %\todo{Use $\neq_4$ to show \#P-hardness.}
    
    \item $\widehat{\mathfrak{m}}_{ij}\widehat f$ has support of size $8$. 
    We can rename the four variables of $\widehat{\mathfrak{m}}_{ij}\widehat f$ in a cyclic permutation, and replace $\braket{00\mid 11}$ and $\braket{11\mid 00}$ by $|\braket{00\mid 11}|$ by \cref{thm: eight-vertex model}.
    We use $\widehat{g}$ to denote this signature.
    Then $M(\widehat g)= \left[\begin{smallmatrix}
    a & 0 & 0 & b\\
    0 & c & d & 0\\
    0 & d & c & 0\\
    b' & 0 & 0 & a\\
    \end{smallmatrix}\right]$ where $a=|\braket{00\mid 11}|, c=\braket{00\mid 00}=\braket{11\mid 11},d=\braket{10\mid 10}=\braket{01\mid 01}$ are positive real numbers and $ b=\braket{10\mid 01}, b'=\braket{01\mid 10}$ are nonzero complex numbers. 
    By \cref{thm: eight-vertex model}, we consider the following three cases:
    (1) $\holant{\neq_2}{\widehat{g}}$ is \#P-hard.
    Then $\holant{\neq_2}{\hF}$ is \#P-hard, we are done.
    (2) There exists $T\in\mathrm{GL}_2(\mathbb{C})$ such that $\#\mathrm{CSP}_2(T\hF)\le_T\holant{\neq_2}{\widehat{g},\hF}.$
    Since $\mathcal{F}$ doesn't satisfy \rm{(\ref{cond: tractable classes})}, $T\hF$ doesn't lie in $\mathscr{A}$ or $\alpha\mathscr{A}$ or $\mathscr{P}$ or $\mathscr{L}$.
    Thus, by \cref{thm: CSP2}, $\#\mathrm{CSP}_2(T\hF)$ is \#P-hard.
    It follows that $\holant{\neq_2}{\hF}$ is \#P-hard.
    (3) $b'= b$ and $a=|b|=c=d$.
    Consider the signature $\widehat{\mathfrak{m}}_{12}\widehat{g}$ realized by mating $\widehat{g}$ and $\widehat{\overline{g}}$ using $\neq_2$.
    $$M(\widehat{\mathfrak{m}}_{12}\widehat{g})=M(\widehat g)N_4M^{\dagger}(\widehat g)= \left[\begin{matrix}
    ab+a\overline{b} & 0 & 0 & a^2+|b|^2\\
    0 &  2cd & c^2+d^2 & 0\\
    0 & c^2+d^2 & 2cd  & 0\\
    a^2+|b|^2 & 0 & 0 &  ab+a\overline{b}\\
    \end{matrix}\right],$$
    which is a matrix with real entries. 
    Moreover, all entries except $ab+a\overline{b}$ are strictly positive.
    If $ab+a\overline{b}=0$, then $\holant{\neq_2}{\widehat{\mathfrak{m}}_{12}\widehat{g}}$ is \#P-hard by the complexity classification of the six-vertex model \cite{six-vertex_model}.
    Otherwise, by \cref{thm: eight-vertex model}, either (1) $\holant{\neq_2}{\widehat{\mathfrak{m}}_{12}\widehat{g}}$ is \#P-hard, and thus $\holant{\neq_2}{\hF}$ is \#P-hard;
    or (2) There exists $T\in\mathrm{GL}_2(\mathbb{C})$ such that $\#\mathrm{CSP}_2(T\hF)\le_T\holant{\neq_2}{\widehat{f},\hF}.$
    Similar to the above argument, $\holant{\neq_2}{\hF}$ is \#P-hard.
    (3) We can realize $\widehat{h}$ with matrix $M(\widehat{h})$=$\left[\begin{smallmatrix}
        1 & 0 & 0 & 1\\
        0 & 1 & 1 & 0\\
        0 & 1 & 1 & 0\\
        1 & 0 & 0 & 1\\
    \end{smallmatrix}\right]$.
    A direct computation shows that $K^{\otimes 2}M(\widehat{h}){(K^{\tt T})}^{\otimes 2}=2\left[\begin{smallmatrix}
        1 & 0 & 0 & 0\\
        0 & 0 & 0 & 0\\
        0 & 0 & 0 & 0\\
        0 & 0 & 0 & 1\\
    \end{smallmatrix}\right].$
    So $\holant{\neq_2}{\widehat{h},\hF}\equiv_T\holant{(\neq_2)K^{-1}}{K\widehat{h},K\hF}\equiv_T\holant{=_2}{=_4,\mathcal{F}}$ by holographic transformation.
    By \cref{lm: CSP2 to holant =4}, $\#\mathrm{CSP}_2(\mathcal{F})\leqslant_T\Holant(=_4, \mathcal F)\equiv_T\holant{\neq_2}{\hF}$.
    Since $\mathcal{F}$ doesn't satisfy condition \rm{(\ref{cond: tractable classes})}, $\#\mathrm{CSP}_2(\mathcal{F})$ is \#P-hard by \cref{thm: CSP2}, and thus $\holant{\neq_2}{\hF}$ is \#P-hard.
    
       \item $\widehat{\mathfrak{m}}_{ij}\widehat f$ has support size $4$. 
       By Cauchy-Schwarz inequality, $M(\widehat{\mathfrak{m}}_{ij}\widehat f)$ is of the form 
    $\left[\begin{smallmatrix}
    a & 0 & 0 & c\\
    0 & 0 & 0 & 0\\
    0 & 0 & 0 & 0\\
    c & 0 & 0 &  \overline{a}\\
    \end{smallmatrix}\right]$
    or  $\left[\begin{smallmatrix}
    0 & 0 & 0 & 0\\
    0 & e & d & 0\\
    0 & d & \overline{e} & 0\\
    0 & 0 & 0 & 0\\
    \end{smallmatrix}\right]$  or $ \left[\begin{smallmatrix}
    0 & 0 & 0 & c\\
    0 & 0 & d & 0\\
    0 & d & 0 & 0\\
    c & 0 & 0 & 0\\
    \end{smallmatrix}\right]$, where $\braket{00\mid 11}=a$, $\braket{00\mid 00}=\braket{11\mid 11}=c$, $\braket{01\mid 01}=\braket{10\mid 10}=d$, and $\braket{01\mid 10}=e.$
    In the first case, by Lemma 6.1 in \cite{eight-vertex_model}, $\holant{\neq_2}{\widehat{\mathfrak{m}}_{ij}\widehat f}$ is \#P-hard unless $|a|^2=c^2.$ 
    By Cauchy-Schwartz inequality, there exists some $\lambda\neq0$ such that $\mathbf{f}_{ij}^{11}=\lambda\mathbf{f}_{ij}^{00}$.
    Also notice that $\mathbf{f}_{ij}^{01}=\mathbf{f}_{ij}^{10}=0$.
    So $\widehat{f}=\widehat{b'}(x_i,x_j)\otimes f_{ij}^{00}$, where $\widehat{b'}(x_i,x_j)$ is the binary signature with matrix
    $\left[\begin{smallmatrix}
        1 & 0\\
        0 & \lambda
    \end{smallmatrix}\right]$.
    Since $|\mathbf{f}_{ij}^{11}|^2=|\mathbf{f}_{ij}^{00}|^2=c$, $|\lambda|=1$.
    Suppose $\lambda=e^{2\mathfrak{i}\theta}$. After a normalization, we have $\widehat{b}(x_i,x_j)\mid \widehat{f}$, where $\widehat{b}(x_i,x_j)\in\widehat{\mathcal{O}}$ is the binary signature with matrix
    $\left[\begin{smallmatrix}
        e^{-\mathfrak{i}\theta} & 0\\
        0 & e^{\mathfrak{i}\theta}
    \end{smallmatrix}\right]$.
    In the second case, by Theorem 3.1 and Corollary 2.6 in \cite{six-vertex_model}, $\holant{\neq_2}{\widehat{\mathfrak{m}}_{ij}\widehat f}$ is \#P-hard unless $d^2=|e|^2$.
    Similar to case one, we can show $\widehat{b}(x_i,x_j)\mid \widehat{f}$, where $\widehat{b}(x_i,x_j)\in\widehat{\mathcal{O}}$.
    In the third case, still by Theorem 3.1 and Corollary 2.6 in \cite{six-vertex_model}, $\holant{\neq_2}{\widehat{\mathfrak{m}}_{ij}\widehat f}$ is \#P-hard unless $c=d$, in which case $M(\mathfrak{m}_{ij}\widehat{f})=\lambda_{ij}N_4$ for $\lambda_{ij}=c$, then we are done.
    \end{itemize}
    To conclude, if there exists a pair of indices $\{i,j\}$ such that $\widehat{\mathfrak{m}}_{ij}\widehat f$ has support size 2 or 8, then $\Holant(\neq_2 \mid \widehat{\mathfrak{m}}_{ij}\widehat f)$ is \#P-hard and thus $\holant{\neq_2}{\hF}$ is \#P-hard.
    Otherwise, for every pair of indices $\{i,j\}$, $\widehat{\mathfrak{m}}_{ij}\widehat f$ has support size 4. 
    Then there exists a nonzero binary signature $\widehat{b}_{ij}\in \widehat{\mathcal{O}}$ such that $\widehat{b}_{ij}(x_i, x_j)\mid \widehat{f}$ or $M(\widehat{\mathfrak{m}}_{ij}\widehat{f})= \lambda_{ij} N_4$ for some real $\lambda_{ij}\neq 0$.  
    The lemma is proved.
\end{proof}

\begin{remark}
    We give a restatement of \cref{lm: pre not 2nd-orth is hard} in the setting of $\Holant(\mathcal{F})$.
    Let $f\in \mathcal{F}$ be a signature of arity $2n\geqslant4$. Then, $\Holant(\mathcal{F})$ is \#P-hard, or $f$ is reducible, or $M(\mathfrak{m}_{ij}f)=\lambda_{ij}I_4$ for some real $\lambda_{ij}\neq 0$.
\end{remark}


\fi


\begin{lemma}\label{lm: not 2nd-orth is hard}
    Let $f\in {\mathcal{F}}$ be a irreducible signature of arity $2n \geqslant 4$.
    If $f$ does not satisfy {\sc 2nd-Orth}, then $\Holant({\mathcal{F}})$ is \#P-hard.
\end{lemma}

\begin{proof}
    Fix an arbitrary pair of indices $\{i,j\}$, we can realize the arity 4 signature $\mathfrak{m}_{ij}f$ by mating.
    By \cref{thm: single quaternary dichotomy}, either $\Holant(\mathfrak{m}_{ij}f)\le_T\Holant(\mathcal{F})$ is \#P-hard, or $\mathfrak{m}_{ij}f$ falls in the six tractable classes of \cref{thm: single quaternary dichotomy}.
    
    The first two classes $\braket{K\mathcal{M}}$ and $\braket{KX\mathcal{M}}$ cannot occur by \cref{lm: 2nd-Orth M and XM cannot happen}.
    
    If $\mathfrak{m}_{ij}f$ is $\mathscr{A}$- or $\mathscr{P}$- or $\mathscr{L}$-transformable, then by \cref{lm: mating is APL-trans implies 2nd-orth}, either $f$ is reducible, contradiction; or $\Holant(\mathcal{F})$ is \#P-hard; or $M(\mathfrak{m}_{ij}f)=\lambda_{ij}I_4$ for some $\lambda_{ij}>0.$
    By \cref{lm: 2nd-Orth uniform}, all $\lambda_{ij}$ have the same value $\lambda$. 
    So $f$ satisfies {\sc 2nd-Orth}, unless $\Holant(\mathcal{F})$ is \#P-hard.

    If $f\in \mathscr{H}$, the lemma follows from \cref{lm: 2nd-orth H type} and \cref{lm: 2nd-Orth uniform}.
    If $f\in \braket{\mathscr{T}}$, the lemma follows from from \cref{lm: 2nd-orth mating reducible} and \cref{lm: 2nd-Orth uniform}.
\end{proof}




